% compendium/ml-alternative-data/execution-control.tex
% Entry migrated on 2026-10-15 from the legacy reinforcement-learning-and-ML-in-execution file (slice 16),
% owned by reinforcement-learning-trading; dispositions and source checks in notes/migration.md, the blocks
% verbatim in notes/artefacts/reinforcement-learning-trading/derivations.tex.

\section{Optimal execution as stochastic control: the benchmark for learned policies}

\begin{qfmodel}[Optimal execution with order flow as stochastic control: the Cartea--Jaimungal objective, dynamics and dynamic-programming equation, as presented by Lehalle and Laruelle]\label{mod:cartea-jaimungal-execution}
\qffield{Statement} A trader sells an inventory $Q_0$ by time $T$, choosing a trading rate $\nu_t$ \citep[app.~A, sec.~A.7.3, pp.~281--282, page images]{alpha-laruelle}. The state is the midprice $S_t$, the net order flow $\mu_t$ of all other traders, the remaining inventory $Q_t$ and the cash account $X_t$, with
\[ dS_t=b\,(\mu_t-\nu_t)\,dt+\sigma\,dW_t,\qquad \hat S_t=S_t-\Bigl(\frac\psi2+\kappa\nu_t\Bigr),\qquad dQ_t=-\nu_t\,dt,\qquad dX_t=\hat S_t\,\nu_t\,dt \]
(p.~282): permanent impact linear in the trading rate, and temporary impact $\kappa\nu_t$ with full immediate decay (p.~281). The reward of a given policy $\nu$ is
\[ H_t(\nu)=\mathbb E_t\Bigl[X_T^\nu+Q_T^\nu\bigl(S_T-A\,Q_T^\nu\bigr)-\phi\int_t^T(Q_s^\nu)^2\,ds\Bigr] \]
(p.~281). The value function is the best such reward, $H^*=\sup_\nu H(\nu)$. By dynamic programming it solves (Interpretation: the book's Proposition~6, p.~282, with the corrections listed under Pitfalls)
\[ 0=\Bigl(\partial_t+\tfrac12\sigma^2\partial_{SS}+\mathcal L^\mu\Bigr)H^*-\phi q^2+\sup_\nu\Bigl\{\nu\Bigl(S-\frac\psi2-\kappa\nu\Bigr)\partial_xH^*+b(\mu-\nu)\,\partial_SH^*-\nu\,\partial_qH^*\Bigr\}, \]
with $H^*(T,x,S,\mu,q)=x+q(S-\psi/2)-Aq^2$, where $\mathcal L^\mu$ is the generator of the flow. Writing $H^*=x+q(S-\psi/2)+h(t,\mu,q)$ gives
\[ \partial_th+\mathcal L^\mu h+b\mu q-\phi q^2+\sup_\nu\bigl\{-\kappa\nu^2-(bq+\partial_qh)\nu\bigr\}=0,\qquad \nu^*=-\frac{bq+\partial_qh}{2\kappa} \]
(p.~283, page image). With $h=h_0+qh_1+q^2h_2$, the coefficient $h_2$ solves a Riccati equation that does not involve $\mu$, $h_1$ carries the expected future flow, and $h_0$ follows by integration (pp.~283--284, page images). The authors' point is that the trader can act on a view of future flows (p.~284). Interpretation, derived and checked numerically: with $b+2h_2=-2\xi$, the optimal rate is
\[ \nu_t^*=\frac{\xi(t)}{\kappa}\,Q_t-\frac{h_1(t,\mu_t)}{2\kappa},\qquad h_1(t,\mu)=b\int_t^Te^{-\frac1\kappa\int_t^s\xi(u)\,du}\,\mathbb E_{t,\mu}[\mu_s]\,ds, \]
so the trader sells in proportion to inventory and sells more slowly when other traders are expected to keep buying.

\qffield{Notation} Units, for a sale: $Q$ in shares; $\nu$ and $\mu$ in shares per unit time, $\nu>0$ selling and $\mu>0$ net buying by others; $S$, $\psi$ and $\hat S$ in currency per share; $b$ in currency per share traded; $\kappa$ in currency $\times$ time per share$^2$; $\phi$ in currency per share$^2$ per unit time; $A$ in currency per share$^2$. Signs: the trader's own selling lowers $S$ permanently by $b$ per share; each fill is $\psi/2+\kappa\nu$ below the midprice; others' net buying raises $S$. $H_t(\nu)$ is the reward of one fixed policy; the legacy notes called it the value function, which is $H^*$.

\qffield{Conditions} The trader observes $S$, $\mu$, $Q$ and $X$ at each time. The flow $\mu$ is Markov with a known generator and does not respond to the trader. The parameters $b$, $\kappa$, $\sigma$, $\psi$, $\phi$ and $A$ are constant: the book writes $b_t$ in the price dynamics but a constant $b$ in the equation. The control is unconstrained: nothing forces $\nu\ge0$ or $Q_t\ge0$. A finite $A$ charges the inventory left at $T$, which is then sold at $S_T-\psi/2-AQ_T$. That is a quadratic penalty, not mandatory completion; completion is forced only in the limit $A\to\infty$. The running penalty $\phi\int Q^2$, which the book calls risk aversion (p.~282), is not in general the variance of implementation shortfall. Interpretation: for a deterministic schedule the price-risk part of the shortfall has variance $\sigma^2\int Q_s^2\,ds$, so $\phi=\lambda\sigma^2$ reproduces a mean--variance penalty. For a policy that adapts to the flow, $\phi$ is a penalty on holding inventory, not a variance. With $\mu\equiv0$ and $A\to\infty$ the optimal holding is $Q_t=Q_0\sinh(\gamma(T-t))/\sinh(\gamma T)$ with $\gamma=\sqrt{\phi/\kappa}$. This is the continuous-time form of the \hyperref[mod:almgren-chriss-execution]{Almgren--Chriss schedule}, with $\kappa$ in the role of $\eta$ and $\phi$ of $\lambda\sigma^2$ (Interpretation; matched numerically to four decimals).

\qffield{Practice} The legacy note's parameter intuition holds under the model's assumptions, with qualifications (Interpretation; own numerical paths with $\mu\equiv0$, $Q_0=1$, $T=1$, base case $\kappa=0.01$, $\phi=0.001$, $A=0.5$, in consistent units):
\begin{itemize}
\item a larger $\phi$ front-loads the sale: tenfold $\phi$ leaves 0.45 rather than 0.50 at $T/2$;
\item a larger $A$ mostly cuts what is left at $T$, 0.2 rather than 1.9 percent for tenfold $A$, and barely changes the early path;
\item a larger $\kappa$ slows and smooths trading: tenfold $\kappa$ leaves 0.58 at $T/2$, and with finite $A$ it also leaves 17 percent unsold at $T$.
\end{itemize}
With flow, the rate also depends on the expected $\mu$. Lehalle and Laruelle suggest Almgren--Chriss when price formation is well understood at a 5--30-minute scale, and this framework when one knows a continuous-time version of price formation focused on trading flows (p.~284).

As the control-theoretic benchmark for the pathway's learned policies (legacy block \emph{Control vs Learning}): stochastic control solves for a policy given a specified model. Simple models are tractable but may oversimplify the market; realistic ones are often intractable \citep[sec.~1, p.~2]{wide-hambly-xu-yang}. Reinforcement learning learns decisions from repeated interaction with a system, without exemplary supervision or a complete model of the environment (sec.~2, p.~3). When the transition law and rewards are unknown the problem becomes one of learning: \emph{model-based} methods learn them explicitly, \emph{model-free} methods implicitly (sec.~2.2, p.~7). The survey motivates RL for execution by Almgren--Chriss's reliance on its dynamics and linear impact, and by its pre-planned schedule, which ignores real-time conditions (sec.~4.2, p.~33, page image). It reports one study that combined the Almgren--Chriss solution with Q-learning, and one double-DQN study that added a final one-second step to force liquidation (p.~34); these studies were not checked. On data, most RL algorithms need more samples than the relevant historical data provide, and financial series are non-stationary, hence the call for sample-efficient methods or good market simulators (sec.~5, p.~49). Interpretation: training against a simulator is interaction, but the simulator must be validated, which is the pathway's checkpoint. Replaying historical order books does not show how the market would have responded to the agent's own orders, so own impact must still come from a model. Learning from logged historical decisions without interaction is a different setting, not covered in the passages read.

\qffield{Pitfalls} Inconsistencies in the appendix (pp.~281--284, page images; corrections are Interpretation, checked algebraically and numerically):
\begin{itemize}
\item Proposition~6 defines the value function maximising the sale as ``$H^*:=\min_\nu H(\nu)$'' while its equation takes a $\sup$; the value is $\sup_\nu H(\nu)$.
\item The terminal condition $x+q(S-\psi/2)-Aq^2$ charges the half-spread on the inventory left at $T$, but the objective's terminal term $Q_T(S_T-AQ_T)$ does not. The ansatz and the $h$-equations follow the terminal condition, so the objective consistent with them is $X_T+Q_T(S_T-\psi/2-AQ_T)-\phi\int Q^2$. Every share then pays $\psi/2$ whether sold before or at $T$, and the spread does not affect the control. With the printed objective, holding shares to $T$ would save $\psi/2$ each: the first-order condition would become $\nu^*=-(bq+\partial_qh+\psi/2)/(2\kappa)$, an incentive to delay.
\item The equation prints $-\partial_qH^*$ inside the supremum. The inventory dynamics $dQ=-\nu\,dt$ give $-\nu\,\partial_qH^*$, and the book's own next step on p.~283 is consistent only with $-\nu\,\partial_qH^*$.
\item Notation: p.~282 says ``add the bid--ask spread $\phi$'' where the spread is $\psi$ and $\phi$ is the running penalty. The equation writes $k$ for $\kappa$ and $b_t$ for $b$, and the terminal conditions for $h_0$ and $h_1$ are printed at $t$ rather than $T$.
\item The half-spread enters $\hat S$ linearly in $\nu$, so a negative rate (buying back) would be credited the half-spread. The model is coherent only for one-directional selling.
\item The printed closed form $h_2=-\xi-b/2$, $\xi(t)=\sqrt{\kappa\phi}\,(1+\zeta e^{2\gamma(T-t)})/(1-\zeta e^{2\gamma(T-t)})$, $\zeta=(A-b/2+\sqrt{\phi\kappa})/(A-b/2-\sqrt{\phi\kappa})$ gives $h_2(T)=A-b$, not $-A$. The Riccati solution has $\xi(t)=\sqrt{\kappa\phi}\,(1+\zeta e^{-2\gamma(T-t)})/(1-\zeta e^{-2\gamma(T-t)})$ with $\zeta=(A-b/2-\sqrt{\kappa\phi})/(A-b/2+\sqrt{\kappa\phi})$; equivalently, the printed expression with the sign of $\xi$ reversed. Checked against numerical integration.
\item The boxed $h_1$ ``when $b=0$'' omits the factor $b$: with $b=0$ the source term vanishes and $h_1\equiv0$, so the flow drops out. The general solution is the $h_1$ above.
\end{itemize}
The legacy notes attribute the model to Cartea and Jaimungal. Their papers and book (\texttt{crypto-cartea}) were not consulted for it; this entry follows Lehalle and Laruelle's account. The legacy taxonomy's ``learn the policy or value function from simulated or historical interaction'' conflates simulator interaction with learning from historical data. Its ``hybrid execution'' (control structure for constraints, machine learning for forecasts, fills and regime detection) is not in the passages read; it is kept in the artefact and stubbed in the pathway's backlog.

\qffield{See also} \hyperref[mod:almgren-chriss-execution]{mean--variance execution}, the discrete deterministic benchmark, owned by the section overview \texttt{execution-trading-costs}; \hyperref[def:price-impact]{price impact} (permanent, temporary and transient impact); \hyperref[def:trading-envelope]{trading envelope} (where a schedule sits in a trading algorithm); \hyperref[def:implementation-shortfall]{implementation shortfall}.

\qffield{Further reading} The legacy blocks it replaces, verbatim, and the numerical checks: \nolinkurl{notes/artefacts/reinforcement-learning-trading/}.

\qffield{Index} optimal execution; stochastic control; Cartea--Jaimungal; order flow; Hamilton--Jacobi--Bellman equation; value function; policy reward; trading rate; permanent impact; temporary impact; terminal inventory penalty; running inventory penalty; Riccati equation; Feynman--Kac; Almgren--Chriss limit; reinforcement learning; model-based RL; model-free RL; market simulator

\qffield{Status} Verification: partly checked 2026-10-15 (against the text layer and page images of a local copy of \texttt{alpha-laruelle}, 2018 World Scientific printing, app.~A, sec.~A.7.3, pp.~281--284, with the HJB reduction, the Riccati and Feynman--Kac solutions and the parameter paths checked numerically; and \texttt{wide-hambly-xu-yang} in the arXiv version 2112.04553v4 of 28 February 2023, not the journal version, pp.~2--4, 7, 32--35 and 49, page images of pp.~32--33). Stability: draft. Currency: model timeless; the survey's account of RL as of February 2023.
\end{qfmodel}
